\documentclass[11pt,a4paper]{article}
\usepackage[margin=1in]{geometry}
\usepackage{amsmath,amssymb,amsthm}
\usepackage{algorithm}
\usepackage{algpseudocode}
\usepackage{booktabs}
\usepackage{hyperref}

\theoremstyle{definition}
\newtheorem{definition}{\textbf{Definition}}[section]
\newtheorem{theorem}{\textbf{Theorem}}[section]
\newtheorem{lemma}[theorem]{\textbf{Lemma}}
\newtheorem{remark}{\textbf{Remark}}[section]

\title{\textbf{Critical Batch Size and Gradient Noise Scale Dynamics in Foundation Model Training}}
\author{\textbf{Team Scaibu} \\ \small Email: \texttt{jaydeep.9664920749@gmail.com} \\ \small \textbf{Architecture and Core Engine Team}}
\date{\textbf{October 2026}}

\begin{document}
\maketitle

\section{Definitions and Cognitive Mental Models}

\subsection{Formal Mathematical Definition}
\textbf{Definition (Gradient Noise Scale and Critical Batch Size $B_{\text{noise}}$):}
During stochastic gradient descent (SGD) optimization of deep neural networks parameterized by $\theta \in \mathbb{R}^P$, the empirical gradient evaluated on an individual data point $x_i \sim \mathcal{D}$ is $g_i \triangleq \nabla_\theta \mathcal{L}(\theta; x_i)$. The true population gradient vector $G \in \mathbb{R}^P$ and parameter covariance matrix $\Sigma \in \mathbb{R}^{P \times P}$ are defined as:
{\boldmath
\begin{equation}
G \triangleq \mathbb{E}_{x \sim \mathcal{D}}\left[ \nabla_\theta \mathcal{L}(\theta; x) \right], \quad \Sigma \triangleq \mathbb{E}_{x \sim \mathcal{D}}\left[ (g_i - G)(g_i - G)^\top \right]
\end{equation}
}
The \textbf{Gradient Noise Scale} (McCandlish et al., 2018), denoted $B_{\text{noise}}$, quantifies the ratio of stochastic gradient variance to the signal power of the mean gradient vector:
{\boldmath
\begin{equation}
B_{\text{noise}} \triangleq \frac{\operatorname{tr}(\Sigma)}{\|G\|_2^2} = \frac{\sum_{j=1}^P \operatorname{Var}(g_{\cdot, j})}{\sum_{j=1}^P G_j^2}
\end{equation}
}
The critical batch size $B_{\text{noise}}$ defines the exact boundary separating two distinct training regimes:
\begin{enumerate}
    \item \textbf{Linear Scaling Regime ($B \ll B_{\text{noise}}$)}: Increasing batch size $B$ reduces stochastic noise proportionally, enabling linear reduction in required training optimization steps with near-100\% compute efficiency.
    \item \textbf{Diminishing Returns Regime ($B \gg B_{\text{noise}}$)}: Increasing $B$ yields negligible variance reduction per token, saturating convergence speed and wasting GPU FLOPs.
\end{enumerate}

\subsection{Expanded Non-Technical Definition (Intuition for Anyone)}
\begin{enumerate}
    \item \textbf{Asking People for Directions}: If you are lost in a city where directions are noisy and confused, asking 10 different people ($B = 10$) gives you a vastly better consensus than asking 1 person. But asking 10,000 people gives you virtually the exact same direction as asking 500 people. The point where asking more people stops helping is the critical batch size $B_{\text{noise}}$.
    \item \textbf{Batch Size Grows During Training}: Early in training, the model knows nothing; the gradients are enormous, clean, and point in a consistent direction ($B_{\text{noise}}$ is small). As the model becomes highly proficient, the true gradient shrinks near a local minimum while random per-sentence variation stays high ($B_{\text{noise}}$ grows from thousands to millions of tokens).
    \item \textbf{Compute Efficiency Frontier}: Training at $B \approx B_{\text{noise}}$ yields the optimal trade-off between wall-clock completion time and total dollar cost of GPU compute.
\end{enumerate}

\subsection{Child-Friendly Mental Model (Physical Metaphor)}
Imagine weighing coins on a windy dock:
\begin{itemize}
    \item \textbf{Small Batch ($B=1$)}: A single gust of wind throws off the scale; you must take hundreds of separate measurements.
    \item \textbf{Optimal Batch ($B = B_{\text{noise}}$)}: Averaging a handful of coins cancels out the wind completely.
    \item \textbf{Wasteful Batch ($B \gg B_{\text{noise}}$)}: Averaging a million coins simultaneously does not eliminate wind any better than 50 coins, but takes a gigantic team to count!
\end{itemize}

\section{Core Mathematical Equations and Definitions}

\subsection{Sample Efficiency and Step Scaling Law}
{\boldmath
\begin{equation}
\mathcal{S}(B) = \mathcal{S}_{\min} \cdot \left( 1 + \frac{B_{\text{noise}}}{B} \right)
\end{equation}
}
\begin{itemize}
    \item \textbf{Formal Definition}: Number of sequential optimization steps $\mathcal{S}(B)$ required to reach a target loss value as a function of batch size $B$.
    \item \textbf{What It Means}: When $B \to \infty$, steps reach the asymptotic minimum $\mathcal{S}_{\min}$. When $B = B_{\text{noise}}$, steps are exactly $2 \times \mathcal{S}_{\min}$.
    \item \textbf{Why It Is Important}: Governs cluster dimensioning, pipeline parallelism schedules, and distributed training schedules across compute topologies.
\end{itemize}

\subsection{Compute Efficiency Metric}
{\boldmath
\begin{equation}
\mathcal{E}(B) \triangleq \frac{\mathcal{E}_{\text{ideal}}}{\mathcal{E}_{\text{actual}}} = \frac{1}{1 + \frac{B}{B_{\text{noise}}}}
\end{equation}
}
\begin{itemize}
    \item \textbf{Formal Definition}: Mathematical ratio between tokens processed and effective optimization progress.
    \item \textbf{What It Means}: At $B = B_{\text{noise}}$, compute efficiency is exactly 50\%. At $B \ll B_{\text{noise}}$, efficiency approaches 100\%.
    \item \textbf{Why It Is Important}: Prevents engineering teams from setting arbitrary distributed batch sizes that burn cluster energy without accelerating model convergence.
\end{itemize}

\section{Rigorous Mathematical Analysis Checklist}

\begin{itemize}
    \item \textbf{Notation}: $B$ (batch size), $P$ (model parameters), $g_i \in \mathbb{R}^P$ (per-sample gradient), $G \in \mathbb{R}^P$ (mean gradient), $\Sigma$ (covariance matrix), $B_{\text{noise}}$ (critical noise scale).
    \item \textbf{Epistemic Status}: Experimentally verified by OpenAI (McCandlish et al., 2018; Kaplan et al., 2020) and Google DeepMind.
    \item \textbf{Dependency Graph}: Per-Sample Gradients $\to$ Sample Mean $G$ \& Covariance Trace $\operatorname{tr}(\Sigma) \to B_{\text{noise}} \to$ Batch Schedule Ramp.
    \item \textbf{Dimensions}: Gradients in $\mathbb{R}^P$; scalar noise scale $B_{\text{noise}} \in \mathbb{R}_+$.
    \item \textbf{Regimes}: Applicable from small scale ($B=32$) to extreme frontier clusters ($B = 4\times 10^6$ tokens).
    \item \textbf{Invariants}: Strict positivity of $B_{\text{noise}} > 0$ whenever stochastic gradient variance $\operatorname{tr}(\Sigma) > 0$.
    \item \textbf{Isomorphisms}: Directly isomorphic to Signal-to-Noise Ratio (SNR) in statistical signal processing and Monte Carlo sample variance estimation.
\end{itemize}

\section{Theoretical Theorems and Formal Proofs}

\begin{theorem}[Unbiased Estimation of Covariance Trace from Mini-Batch Gradients]
Let $\{g_1, \dots, g_B\}$ be $B \ge 2$ independent and identically distributed per-sample gradient vectors drawn from distribution $\mathcal{D}$ with mean $G = \mathbb{E}[g_i]$ and covariance $\Sigma = \operatorname{Cov}(g_i)$. Define the sample mean $\bar{g} = \frac{1}{B} \sum_{i=1}^B g_i$ and the sample variance estimator for coordinate $j \in \{1, \dots, P\}$:
\begin{equation}
s_j^2 = \frac{1}{B - 1} \sum_{i=1}^B \left( g_{i, j} - \bar{g}_j \right)^2
\end{equation}
Then the sample trace $\operatorname{tr}(\hat{\Sigma}) = \sum_{j=1}^P s_j^2$ is an unbiased estimator of the population trace $\operatorname{tr}(\Sigma)$:
{\boldmath
\begin{equation}
\mathbb{E}\left[ \operatorname{tr}(\hat{\Sigma}) \right] = \operatorname{tr}(\Sigma)
\end{equation}
}
\end{theorem}

\begin{proof}
For each coordinate $j \in \{1, \dots, P\}$, let $X_i = g_{i, j}$. By standard sample variance theory:
\begin{equation}
\mathbb{E}\left[ s_j^2 \right] = \mathbb{E}\left[ \frac{1}{B - 1} \sum_{i=1}^B (X_i - \bar{X})^2 \right]
\end{equation}
Expanding the sum of squares:
\begin{equation}
\sum_{i=1}^B (X_i - \bar{X})^2 = \sum_{i=1}^B X_i^2 - B \bar{X}^2
\end{equation}
Taking expectations:
\begin{equation}
\mathbb{E}\left[ \sum_{i=1}^B X_i^2 \right] = B \left( \operatorname{Var}(X) + (\mathbb{E}[X])^2 \right) = B \Sigma_{jj} + B G_j^2
\end{equation}
For the sample mean $\bar{X}$:
\begin{equation}
\mathbb{E}\left[ \bar{X}^2 \right] = \operatorname{Var}(\bar{X}) + (\mathbb{E}[\bar{X}])^2 = \frac{\Sigma_{jj}}{B} + G_j^2
\end{equation}
Thus:
\begin{equation}
\mathbb{E}\left[ \sum_{i=1}^B (X_i - \bar{X})^2 \right] = B \Sigma_{jj} + B G_j^2 - B\left( \frac{\Sigma_{jj}}{B} + G_j^2 \right) = (B - 1) \Sigma_{jj}
\end{equation}
Dividing by $B - 1$ yields $\mathbb{E}[s_j^2] = \Sigma_{jj}$.
Summing over all dimensions $j = 1, \dots, P$:
\begin{equation}
\mathbb{E}\left[ \operatorname{tr}(\hat{\Sigma}) \right] = \sum_{j=1}^P \mathbb{E}[s_j^2] = \sum_{j=1}^P \Sigma_{jj} = \operatorname{tr}(\Sigma)
\end{equation}
which completes the proof.
\end{proof}

\begin{theorem}[Pareto Optimality of Critical Batch Size Schedule]
Under a quadratic loss landscape $\mathcal{L}(\theta) = \frac{1}{2} \theta^\top H \theta$ with learning rate $\eta(B) = \frac{\eta_0}{1 + B_{\text{noise}} / B}$, the step progress per unit of wall-clock time $\mathcal{R}_{\text{time}}(B)$ and per unit of total compute tokens $\mathcal{R}_{\text{FLOPs}}(B)$ cannot be simultaneously increased beyond $B = B_{\text{noise}}$. Specifically, the Pareto frontier satisfies:
{\boldmath
\begin{equation}
\frac{d \mathcal{S}}{d B} \cdot \frac{B}{\mathcal{S}} = -\frac{B_{\text{noise}}}{B + B_{\text{noise}}}
\end{equation}
}
which reaches elasticity $-0.5$ at $B = B_{\text{noise}}$.
\end{theorem}

\begin{proof}
From the step scaling relation $\mathcal{S}(B) = \mathcal{S}_{\min} \left(1 + \frac{B_{\text{noise}}}{B}\right)$, compute the derivative with respect to $B$:
\begin{equation}
\frac{d \mathcal{S}}{d B} = -\mathcal{S}_{\min} \frac{B_{\text{noise}}}{B^2}
\end{equation}
The elasticity of steps with respect to batch size is:
\begin{equation}
\varepsilon_{\mathcal{S}, B} \triangleq \frac{d \mathcal{S}}{d B} \frac{B}{\mathcal{S}(B)} = \frac{-\mathcal{S}_{\min} \frac{B_{\text{noise}}}{B^2} \cdot B}{\mathcal{S}_{\min} \left( 1 + \frac{B_{\text{noise}}}{B} \right)} = \frac{-\frac{B_{\text{noise}}}{B}}{\frac{B + B_{\text{noise}}}{B}} = -\frac{B_{\text{noise}}}{B + B_{\text{noise}}}
\end{equation}
When $B \ll B_{\text{noise}}$, $\varepsilon_{\mathcal{S}, B} \approx -1$, meaning a 1\% increase in batch size reduces steps by 1\% (perfect speedup).
When $B \gg B_{\text{noise}}$, $\varepsilon_{\mathcal{S}, B} \to 0$, meaning batch size increases yield zero step reduction.
At $B = B_{\text{noise}}$, $\varepsilon_{\mathcal{S}, B} = -\frac{1}{2}$, establishing the exact inflection point of diminishing marginal returns.
\end{proof}

\section{Architectural Questions and Real-World Trade-Offs}

\subsection{Architectural Question 1: Batch Size Ramp Schedules}
\textbf{Question:} Why do modern LLM training recipes start with small batch sizes and ramp up over time?\\
\textbf{Answer:} Gradient noise scale $B_{\text{noise}}$ is non-stationary; it is small during early training when gradient signal $\|G\|^2$ is massive, and increases by orders of magnitude as training converges. Ramping batch size (e.g., from 4M tokens to 16M tokens over the first 20\% of training) maintains near-optimal compute efficiency throughout the run, saving hundreds of thousands of GPU-hours compared to fixed-batch execution.

\subsection{Architectural Question 2: Per-Sample Gradient Memory Overhead}
\textbf{Question:} How is $B_{\text{noise}}$ estimated in production without materializing per-sample gradients in GPU VRAM?\\
\textbf{Answer:} Storing $B$ individual gradient vectors of size $P$ would exhaust HBM memory ($O(B \cdot P)$ bytes). Production runners estimate $B_{\text{noise}}$ by evaluating two different sub-batch gradient norms: $G_{\text{small}}$ (batch $B_1$) and $G_{\text{large}}$ (batch $B_2$). By leveraging the relation $\mathbb{E}[\|G_B\|^2] = \|G\|^2 + \frac{\operatorname{tr}(\Sigma)}{B}$, the noise scale is algebraically extracted using only 2 scalar norms without per-sample gradient storage.

\section{Algorithmic Implementation Specification}

\begin{algorithm}[H]
\caption{Critical Batch Size and Gradient Noise Scale Estimation}
\begin{algorithmic}[1]
\Require Per-sample gradients $\mathbf{g} \in \mathbb{R}^{B_{\text{sample}} \times P}$, step $t$, total steps $T$, bounds $[B_{\min}, B_{\max}]$
\Ensure Noise scale $B_{\text{noise}}$, covariance trace $\operatorname{tr}(\Sigma)$, gradient norm $\|G\|^2$, recommended batch size $B^*$, efficiency $\mathcal{E}$
\State $B_{\text{sample}} \gets \operatorname{len}(\mathbf{g}), \quad P \gets \operatorname{len}(\mathbf{g}[0])$
\State Compute sample mean gradient: $\bar{\mathbf{G}} \gets \frac{1}{B_{\text{sample}}} \sum_{i=1}^{B_{\text{sample}}} \mathbf{g}_i$
\State Compute squared mean norm: $\|\bar{\mathbf{G}}\|^2 \gets \sum_{j=1}^P \bar{G}_j^2$
\State Initialize trace sum $\operatorname{tr}(\Sigma) \gets 0.0$
\For{$j = 0$ \textbf{to} $P - 1$}
    \State $\operatorname{Var}_j \gets \frac{1}{B_{\text{sample}} - 1} \sum_{i=1}^{B_{\text{sample}}} (g_{i, j} - \bar{G}_j)^2$
    \State $\operatorname{tr}(\Sigma) \gets \operatorname{tr}(\Sigma) + \operatorname{Var}_j$
\EndFor
\State $B_{\text{noise}} \gets \operatorname{tr}(\Sigma) / \max(10^{-12}, \|\bar{\mathbf{G}}\|^2)$
\State $\text{progress} \gets \min(1.0, \max(0.0, t / T))$
\State $B^* \gets \operatorname{clip}\left( \operatorname{round}\left( B_{\text{noise}} \cdot (0.5 + 0.5 \cdot \text{progress}) \right), B_{\min}, B_{\max} \right)$
\State $\mathcal{E} \gets 1.0 / (1.0 + B^* / \max(1.0, B_{\text{noise}}))$
\State \Return $(B_{\text{noise}}, \operatorname{tr}(\Sigma), \|\bar{\mathbf{G}}\|^2, B^*, \mathcal{E})$
\end{algorithmic}
\end{algorithm}

\section{Big-$\Theta$ Complexity Analysis}

\begin{itemize}
    \item \textbf{Mean Gradient Reduction}: $\Theta(B_{\text{sample}} \cdot P)$ vector accumulation passes.
    \item \textbf{Coordinate Variance Accumulation}: $\Theta(B_{\text{sample}} \cdot P)$ element-wise differences.
    \item \textbf{Norm and Ratio Computation}: $\Theta(P)$ linear vector pass.
    \item \textbf{Total Time Complexity}: $\Theta(B_{\text{sample}} \cdot P)$ strictly optimal with respect to input tensor size.
    \item \textbf{Space Complexity}: $\Theta(P)$ auxiliary buffer for mean gradient storage.
\end{itemize}

\section{Single Canonical Repository Reference}

The complete deterministic scalar implementation, verification suite, and unit test benchmarks are published at:
\begin{center}
\url{https://github.com/Chief-Strategist-J/llm-obs-infra}
\end{center}

\end{document}
